\chapter{Optimization and resampling}
\label{chap:optimres}
%=============================================================================

Chapter~\ref{chap:vmc} left two loose ends.  The variational parameters were
located by scanning a grid, which worked because there were only two of them;
and the error bar was obtained from the autocorrelation time, which is
serviceable but not what a production calculation uses.  Both are addressed
here, and the two are more closely related than they look: the first is about
spending computer time cheaply while hunting for a minimum, the second about
spending it lavishly once the minimum is found and then reporting honestly
what was achieved.

The structure of a real calculation has three stages.

\begin{svgraybox}
\begin{enumerate}
\item \textbf{Optimise.}  Compute the gradient of the energy with respect to
  the variational parameters -- which turns out to be another Monte Carlo
  average, costing nothing extra -- and follow it downhill.  Short, noisy runs
  suffice: we are locating a minimum, not measuring an energy.
\item \textbf{Produce.}  With the parameters fixed, run once with as many
  samples as the budget allows.  Millions, and on a cluster, billions.
\item \textbf{Resample.}  The samples are correlated, so
  $\sigma/\sqrt{N}$ is wrong.  Blocking, the bootstrap or the jackknife give
  the honest error.
\end{enumerate}
\end{svgraybox}

\noindent
The programs are in \texttt{vmcoptimise.py}, which imports the trial function,
local energy and quantum force of Chapter~\ref{chap:vmc} unchanged.

%----------------------------------------------------------------------
\section{Why a grid does not scale}
\label{sec:optgrid}

Table~\ref{tab:13-surface} was a $5\times5$ grid: twenty-five Monte Carlo runs
to locate two parameters, and the answer only as good as the mesh.  With $p$
parameters and $m$ points along each, the cost is $m^p$ runs.  For $p=2$ and
$m=5$ that is twenty-five; for $p=10$ it is ten million; and for the
neural-network trial functions of the later chapters, where $p$ runs to
thousands, it is not a number worth writing down.

Worse, a grid throws away information.  Each run already knows which way is
downhill -- the samples that produced the energy can produce its derivative
too -- and a scan discards that and starts again at the next node.  What we
need is the gradient.

%----------------------------------------------------------------------
\section{The gradient of the energy}
\label{sec:optgradient}

\index{energy gradient}\index{variational Monte Carlo!optimisation}
Let $\theta$ be a variational parameter and write
\begin{equation}
  \bar{E}_{\theta} = \frac{d\mean{E_L[\theta]}}{d\theta},
  \qquad
  \bar{\psi}_{\theta} = \frac{d\psi[\theta]}{d\theta} .
  \label{eq:14-definitions}
\end{equation}
Differentiating
$E=\bracket{\psi}{\hat{H}}{\psi}/\braket{\psi}{\psi}$ and using the
hermiticity of $\hat{H}$ gives, after a short calculation set as
Exercise~\ref{ex:14-gradient},
\begin{equation}
  \boxed{\;
  \bar{E}_{\theta}
   = 2\left(
     \meanb{\frac{\bar{\psi}_{\theta}}{\psi}\,E_L}
     - \meanb{\frac{\bar{\psi}_{\theta}}{\psi}}\meanb{E_L}
     \right) \;}
  \label{eq:14-gradient}
\end{equation}
This is the central formula of the chapter, and three things about it are
worth noticing.

First, it is a \emph{covariance}: the gradient is twice the covariance of the
local energy with the logarithmic derivative of the trial function,
\begin{equation}
  O_{\theta}(\bm{R}) = \frac{\bar{\psi}_{\theta}(\bm{R})}{\psi(\bm{R})}
   = \frac{\partial\ln\psi(\bm{R})}{\partial\theta} .
  \label{eq:14-Otheta}
\end{equation}
If the local energy and the logarithmic derivative are uncorrelated the
gradient vanishes, which is another way of saying that at the minimum the
trial function has stopped being able to tell good configurations from bad
ones.

Second, it needs only \emph{first} derivatives of $\ln\psi$ with respect to
the parameters -- no second derivatives, no derivatives of the local energy,
and no extra sampling.  The same walk that measures the energy accumulates
$\mean{O_\theta}$ and $\mean{O_\theta E_L}$ in the same loop.

Third, and least obviously, it is much better behaved than a finite
difference.  Estimating $\partial E/\partial\theta$ by running at
$\theta\pm h$ and subtracting requires two runs per parameter, and the
subtraction cancels most of the signal while the statistical errors add.
Equation~(\ref{eq:14-gradient}) has no cancellation of that kind.

\paragraph{For our trial function.}
With $\Psi_T$ of Eq.~(\ref{eq:13-trial}) the two logarithmic derivatives are
elementary,
\begin{equation}
  O_{\alpha} = -\frac{\omega}{2}\left(r_1^2+r_2^2\right),
  \qquad
  O_{\beta} = -\frac{r_{12}^2}{\left(1+\beta r_{12}\right)^2} ,
  \label{eq:14-Odot}
\end{equation}
and \texttt{vmcoptimise.py} checks both against finite differences, agreeing
to $2\times10^{-9}$.  It then checks the assembled gradient
Eq.~(\ref{eq:14-gradient}) against a finite difference of the energy itself;
Table~\ref{tab:14-gradient} shows the comparison.

\begin{table}[h]
\centering
\renewcommand{\arraystretch}{1.25}
\setlength{\tabcolsep}{10pt}
\begin{tabular}{rrrrrr}
\toprule
& & \multicolumn{2}{c}{$\partial E/\partial\alpha$} &
    \multicolumn{2}{c}{$\partial E/\partial\beta$}\\
\cmidrule(lr){3-4}\cmidrule(l){5-6}
$\alpha$ & $\beta$ & Eq.~(\ref{eq:14-gradient}) & finite difference
 & Eq.~(\ref{eq:14-gradient}) & finite difference\\
\midrule
$0.90$ & $0.30$ & $-0.41171$ & $-0.42413$ & $-0.29347$ & $-0.29289$\\
$1.00$ & $0.40$ & $\phantom{-}0.02905$ & $\phantom{-}0.03823$ &
  $\phantom{-}0.01227$ & $\phantom{-}0.01853$\\
$1.05$ & $0.50$ & $\phantom{-}0.22099$ & $\phantom{-}0.21838$ &
  $\phantom{-}0.10431$ & $\phantom{-}0.10902$\\
\bottomrule
\end{tabular}
\caption{The analytic gradient against a central difference of the energy with
$h=0.01$, $40\,000$ cycles per run.  They agree within the Monte Carlo noise,
and the analytic route costs one run instead of four.  Computed by
\texttt{vmcoptimise.py}.}
\label{tab:14-gradient}
\end{table}

%----------------------------------------------------------------------
\section{Descending the surface}
\label{sec:optdescent}

\index{gradient descent}\index{steepest descent}
With a gradient in hand the simplest scheme is \emph{steepest descent},
\begin{equation}
  \bm{\theta}_{k+1} = \bm{\theta}_k
   - \eta\,\nabla_{\bm{\theta}}E(\bm{\theta}_k) ,
  \label{eq:14-steepestdescent}
\end{equation}
with $\eta$ the learning rate.  For a convex function and a small enough
$\eta$ this converges; the difficulty is that ``small enough'' depends on the
curvature, and the curvature is exactly what the gradient does not know.  If
the surface is a narrow valley -- the Hessian badly conditioned -- steepest
descent zig-zags across it and creeps along it, and the number of iterations
grows with the condition number.

Table~\ref{tab:14-descent} runs it on the quantum dot, starting deliberately
far from the minimum.  Two features are worth pointing out.  The energy falls
from $3.1128$ to within $2\times10^{-3}$ of the exact answer in about ten
steps; and each step uses only $6000$ Monte Carlo cycles, which would be far
too few to \emph{report} an energy but is ample to decide which way is down.
That asymmetry -- cheap runs while optimising, one expensive run at the end --
is the reason the two halves of this chapter belong together.

\begin{table}[h]
\centering
\renewcommand{\arraystretch}{1.2}
\setlength{\tabcolsep}{13pt}
\begin{tabular}{rllll}
\toprule
iteration & $\alpha$ & $\beta$ & energy & $|\nabla E|$\\
\midrule
$1$  & $0.85000$ & $0.20000$ & $3.112849$ & $1.321241$\\
$2$  & $0.89515$ & $0.24822$ & $3.046038$ & $0.722958$\\
$3$  & $0.92182$ & $0.27262$ & $3.029580$ & $0.507078$\\
$4$  & $0.94113$ & $0.28906$ & $3.013600$ & $0.366792$\\
$6$  & $0.96631$ & $0.31072$ & $3.007260$ & $0.227005$\\
$10$ & $0.99126$ & $0.33328$ & $3.002706$ & $0.099140$\\
$15$ & $1.00300$ & $0.34749$ & $3.000419$ & $0.045287$\\
$20$ & $1.00623$ & $0.35534$ & $3.000180$ & $0.028328$\\
$25$ & $1.00660$ & $0.36082$ & $3.001775$ & $0.017045$\\
\bottomrule
\end{tabular}
\caption{Steepest descent from $(\alpha,\beta)=(0.85,0.20)$ with
$\eta=0.05$ and $6000$ cycles per step.  The gradient is noisy, which is why
the last few energies wander at the level of $10^{-3}$; the parameters are
nonetheless located well enough.  Computed by \texttt{vmcoptimise.py}.}
\label{tab:14-descent}
\end{table}

\paragraph{Doing better than plain descent.}
Three standard improvements, in increasing order of sophistication.

\emph{Momentum} accumulates a velocity,
\begin{equation}
  \bm{v}_{k+1} = \gamma\bm{v}_k + \eta\nabla E,
  \qquad
  \bm{\theta}_{k+1} = \bm{\theta}_k - \bm{v}_{k+1} ,
  \label{eq:14-momentum}
\end{equation}
which damps the zig-zag across a valley and accelerates along it.  With a
\emph{stochastic} gradient it does something else as well, and arguably more
useful: it averages successive gradients, so the noise in each individual
estimate matters less.

\emph{Quasi-Newton} methods build an approximation to the inverse Hessian from
successive gradients and take
$\bm{\theta}_{k+1}=\bm{\theta}_k-\bm{H}_k^{-1}\nabla E$.  Newton's method
itself requires the Hessian, which for a Monte Carlo energy is expensive and
noisy; Broyden's method and the BFGS update instead impose the \emph{secant
condition}
\begin{equation}
  \bm{H}_{k+1}\left(\bm{\theta}_{k+1}-\bm{\theta}_k\right)
   = \nabla E_{k+1} - \nabla E_k
  \label{eq:14-secant}
\end{equation}
and make the smallest change to $\bm{H}$ consistent with it.  These are the
work-horses of deterministic optimisation, and \texttt{scipy.optimize}
provides them; the caveat is that they assume the gradient is exact, and a
stochastic gradient can confuse the curvature estimate badly.

\emph{Stochastic reconfiguration} is the method built for this problem.  The
trouble with Eq.~(\ref{eq:14-steepestdescent}) is that it treats the
parameters as if they lived in a flat Euclidean space, when what we care about
is the distance between the \emph{wave functions} they produce.  Doubling a
parameter that barely affects $\Psi_T$ should cost a large step; changing one
that reshapes it entirely should cost a small one.  The natural metric is the
covariance of the logarithmic derivatives,
\begin{equation}
  S_{kl} = \mean{O_k O_l} - \mean{O_k}\mean{O_l},
  \label{eq:14-metric}
\end{equation}
the \emph{quantum geometric tensor}, which up to a factor of four is the
quantum Fisher information.  Preconditioning with it,
\begin{equation}
  \boxed{\;
  \bm{\theta}_{k+1} = \bm{\theta}_k
   - \eta\left(\bm{S}+\varepsilon\bm{I}\right)^{-1}\nabla_{\bm{\theta}}E \;}
  \label{eq:14-sr}
\end{equation}
makes the step independent of how the parameters happen to be parametrised.
The small shift $\varepsilon$ keeps $\bm{S}$ invertible, which matters because
it is often nearly singular -- two parameters that do almost the same thing to
$\Psi_T$ produce almost identical columns.  Equation~(\ref{eq:14-sr}) is also
imaginary-time evolution projected onto the variational manifold, which is the
route by which it is usually derived, and it is the standard optimiser for
neural-network wave functions.

Table~\ref{tab:14-optimisers} compares the three from the same starting point
with the same budget.

\begin{table}[h]
\centering
\renewcommand{\arraystretch}{1.25}
\setlength{\tabcolsep}{12pt}
\begin{tabular}{lrlll}
\toprule
method & steps & $\alpha$ & $\beta$ & final energy\\
\midrule
gradient descent           & $25$ & $1.00639$ & $0.36165$ & $3.001775$\\
with momentum              & $25$ & $0.99876$ & $0.37802$ & $3.001499$\\
stochastic reconfiguration & $25$ & $0.98977$ & $0.39651$ & $3.001494$\\
\bottomrule
\end{tabular}
\caption{The same start, $(0.85,0.20)$, and the same $6000$ cycles per step.
All three reach the minimum; the differences are in how directly they get
there, and the ordering would be far more pronounced with more parameters or
a worse-conditioned surface.  Computed by \texttt{vmcoptimise.py}.}
\label{tab:14-optimisers}
\end{table}

\begin{notebox}
\textbf{Why optimisation gets a chapter.} For two parameters any of these
methods works and the choice hardly matters.  The reason to set them up
carefully is what comes later: a neural-network trial function has thousands
of parameters, no grid is conceivable, the surface is not convex, and the
gradient is stochastic.  Equation~(\ref{eq:14-gradient}) and
Eq.~(\ref{eq:14-sr}) are exactly the tools that carry over, and the
machine-learning literature knows the first of them as the REINFORCE or
score-function estimator.
\end{notebox}

%----------------------------------------------------------------------
\section{The production run}
\label{sec:optproduction}

Once the parameters are fixed the calculation changes character.  There is
nothing left to decide, only statistics to accumulate, and that parallelises
perfectly: walkers are independent, so an ensemble of them can be advanced
with array operations rather than a Python loop.  With a thousand walkers
taking two thousand steps each, \texttt{vmcoptimise.py} collects two million
local-energy evaluations in about a second, at an acceptance rate of $87\%$.

What it returns is the walker-averaged local energy at each step.  That series
is still autocorrelated \emph{in time} -- averaging over independent walkers
divides the variance and the autocovariance by the same factor and leaves the
normalised autocorrelation function, and therefore $\tau$, unchanged.  The
error analysis of Chapter~\ref{chap:statelements} is still needed.

%----------------------------------------------------------------------
\section{Resampling}
\label{sec:optresampling}

\index{resampling}\index{blocking}\index{bootstrap}\index{jackknife}
Chapter~\ref{chap:statelements} established that a Markov chain gives
correlated samples, that $\sigma_m^2=\tau\sigma^2/n$ rather than $\sigma^2/n$,
and that the naive error is therefore always too small.  It also introduced
blocking.  Here we set out the three standard estimators properly, because a
production calculation must report one of them.

\paragraph{Blocking.}
Replace the data by the averages of neighbouring pairs,
\begin{equation}
  x_i^{(k+1)} = \tfrac12\left(x_{2i-1}^{(k)} + x_{2i}^{(k)}\right),
  \label{eq:14-blocking}
\end{equation}
which halves the number of samples and leaves the mean unchanged.  Each
transformation reduces the correlation between neighbours, so the naive error
computed from the blocked data grows -- until the blocks are longer than the
correlation time, after which it stops growing.  The plateau is the true
error.  Reading the plateau by eye is unsatisfactory, and Jonsson turned it
into a test: the remaining autocovariance after $k$ transformations,
accumulated over the coarser levels, is compared with a $\chi^2$ quantile, and
the first $k$ that passes is used.  That is what \texttt{vmcoptimise.py}
implements, and it makes blocking a method with no free parameters at all.

\paragraph{The bootstrap.}
Draw $B$ new data sets of the same size from the original, \emph{with
replacement}, compute the estimator on each, and take the standard deviation
of the $B$ results.  The essential point for us is that the resampling must be
done in \emph{blocks} rather than one sample at a time: drawing individual
samples destroys the correlations and reproduces the naive error exactly.  The
block length has to be at least a few correlation times, which reintroduces a
parameter that blocking does not need.

\paragraph{The jackknife.}
Divide the data into $n$ blocks, and let $\hat\theta_{(i)}$ be the estimator
computed with block $i$ removed.  Then
\begin{equation}
  \sigma^2_{\mathrm{jack}}
   = \frac{n-1}{n}\sum_{i=1}^{n}
     \left(\hat\theta_{(i)} - \bar{\theta}\right)^2 .
  \label{eq:14-jackknife}
\end{equation}
For the mean this is exact and reduces to the blocked standard error, so it
earns nothing here.  It earns its keep, as does the bootstrap, for
\emph{non-linear} estimators -- a ratio of averages, a variance, a fitted
parameter, an extrapolation -- where there is no formula for the error and
resampling is the only practical route.  The blocking method handles the mean
and nothing else.

\paragraph{What they give.}
Table~\ref{tab:14-resampling} applies all four estimators to two production
runs of the same system: one with the well-tuned time step of
Chapter~\ref{chap:vmc}, and one deliberately badly tuned.

\begin{table}[h]
\centering
\renewcommand{\arraystretch}{1.25}
\setlength{\tabcolsep}{12pt}
\begin{tabular}{lll}
\toprule
estimator & $\Delta t = 0.5$ & $\Delta t = 0.02$\\
\midrule
correlation time $\tau$ & $1.27$ & $13.95$\\
\midrule
naive $\sigma/\sqrt{N}$ & $0.00003276$ & $0.00003377$\\
$\sqrt{\tau}\times$ naive & $0.00003690$ & $0.00012614$\\
blocking & $0.00004858$ & $0.00015405$\\
bootstrap & $0.00003486$ & $0.00010611$\\
jackknife & $0.00003478$ & $0.00010427$\\
\midrule
resampled / naive & $1.48$ & $4.56$\\
\bottomrule
\end{tabular}
\caption{Error estimates for two production runs of two million and 1.6
million samples.  With a well-tuned time step the naive estimate is only
modestly wrong; with a badly tuned one it is wrong by a factor of four, and
silently.  Computed by \texttt{vmcoptimise.py}.}
\label{tab:14-resampling}
\end{table}

Two honest observations.  When the sampling is well tuned, $\tau$ is close to
one and the naive error is not far off -- but there is no way to know that
without measuring.  And the three resampling estimates do not agree exactly:
blocking with Jonsson's stopping rule is deliberately conservative and comes
out largest.  They agree on the order of magnitude, which is what matters, and
all three are right in the case where the naive estimate is badly wrong.

%----------------------------------------------------------------------
\section{The result}
\label{sec:optresult}

Putting the three stages together -- optimise with cheap runs, produce with an
expensive one, resample the error -- gives
\begin{equation}
  E_{\mathrm{VMC}} = 3.000549 \pm 0.000049 ,
  \label{eq:14-final}
\end{equation}
from two million samples at $\alpha=1.00$, $\beta=0.40$.  Set against the
exact $3$ and the results of Chapter~\ref{chap:applications}:

\begin{center}
\renewcommand{\arraystretch}{1.2}
\setlength{\tabcolsep}{14pt}
\begin{tabular}{lll}
\toprule
method & energy & error\\
\midrule
Hartree-Fock, $42$ orbitals & $3.161921$ & $1.6\times10^{-1}$\\
MP2, $42$ orbitals & $3.027038$ & $2.7\times10^{-2}$\\
CCSD, $42$ orbitals & $3.013626$ & $1.4\times10^{-2}$\\
VMC, this chapter & $3.000549(49)$ & $5.5\times10^{-4}$\\
\midrule
exact (Taut) & $3.000000$ & ---\\
\bottomrule
\end{tabular}
\end{center}

\noindent
The remaining $5.5\times10^{-4}$ is eleven standard errors wide, and it is not
statistical: it is what a two-parameter trial function cannot represent.  The
point of this chapter is precisely that the error bar is now small enough, and
trustworthy enough, for that statement to be meaningful.  With the naive error
of $3.3\times10^{-5}$ one would have quoted a result seventeen standard errors
from the exact answer and concluded, wrongly, that something was broken.

%----------------------------------------------------------------------
\section{Summary}
\label{sec:optsummary}

The gradient of a variational energy with respect to its parameters,
Eq.~(\ref{eq:14-gradient}), is a covariance and costs nothing beyond the
energy itself.  With it, the grid scan of Chapter~\ref{chap:vmc} is replaced
by descent -- plain, with momentum, quasi-Newton, or preconditioned by the
metric of the variational manifold -- and the number of parameters stops being
the limiting consideration.  Once they are fixed, the production run is pure
statistics and parallelises trivially, and the error on it must come from
blocking, the bootstrap or the jackknife rather than from $\sigma/\sqrt{N}$.

What remains is the trial function.  Every result in this chapter is bounded
above by what two parameters can express, and no amount of sampling or
optimisation will change that.  Two directions lead out.  Diffusion Monte
Carlo removes the dependence on $\Psi_T$ by projecting onto the ground state
in imaginary time, using the drift and Green's function of
Section~\ref{sec:vmcimportance}.  And if the difficulty is that we cannot
guess a good enough functional form, the alternative is not to guess: to
parametrise $\Psi_T$ by a neural network and let Eq.~(\ref{eq:14-gradient})
and Eq.~(\ref{eq:14-sr}) do the rest.  Both are taken up in the chapters that
follow.

%=============================================================================
\section{Exercises}
%=============================================================================

\subsection*{Exercise 1: the energy gradient}
\label{ex:14-gradient}
\addcontentsline{toc}{subsection}{Exercise 1: the energy gradient}

\begin{enumerate}
  \item[a)] Derive Eq.~(\ref{eq:14-gradient}).
  \item[b)] Find the corresponding expression for the derivative of the
  variance.
  \item[c)] Show that the gradient vanishes when $\Psi_T$ is an exact
  eigenstate, for every parameter, and say why that is obvious in advance.
\end{enumerate}

\paragraph{Answer to a).}
Write $E=\bracket{\psi}{\hat{H}}{\psi}/\braket{\psi}{\psi}$ and differentiate,
\[
  \bar{E}_{\theta}
   = \frac{\bracket{\bar\psi_\theta}{\hat{H}}{\psi}
      + \bracket{\psi}{\hat{H}}{\bar\psi_\theta}}{\braket{\psi}{\psi}}
   - \frac{\bracket{\psi}{\hat{H}}{\psi}}{\braket{\psi}{\psi}}\,
     \frac{\braket{\bar\psi_\theta}{\psi}
      + \braket{\psi}{\bar\psi_\theta}}{\braket{\psi}{\psi}} .
\]
For a real trial function the two terms in each numerator are equal, so both
brackets give a factor of two.  Dividing numerator and denominator by
$\braket{\psi}{\psi}$ turns each ratio into an average over $|\psi|^2$:
$\bracket{\bar\psi_\theta}{\hat{H}}{\psi}/\braket{\psi}{\psi}
= \mean{(\bar\psi_\theta/\psi)E_L}$ by inserting $\psi/\psi$ exactly as in
Eq.~(\ref{eq:13-localenergyaverage}), and
$\braket{\bar\psi_\theta}{\psi}/\braket{\psi}{\psi}
= \mean{\bar\psi_\theta/\psi}$.  Collecting gives
Eq.~(\ref{eq:14-gradient}).

\paragraph{Answer to b).}
With $\sigma^2=\mean{E_L^2}-\mean{E_L}^2$ the same manipulation applied to
$\mean{E_L^2}$ gives
\[
  \frac{d\sigma^2}{d\theta}
   = 2\left(\mean{O_\theta E_L^2} - \mean{O_\theta}\mean{E_L^2}\right)
     + 2\mean{\partial_\theta E_L\,E_L}
     - 2\bar{E}_\theta\mean{E_L} ,
\]
where the middle term appears because $E_L$ itself depends on $\theta$,
unlike in a).  This is why variance minimisation, although attractive in
principle, is more work: it needs $\partial_\theta E_L$, which requires
differentiating the Hamiltonian acting on $\Psi_T$.

\paragraph{Answer to c).}
If $\Psi_T$ is an eigenstate then $E_L$ is a constant, so
$\mean{O_\theta E_L} = E_L\mean{O_\theta} = \mean{O_\theta}\mean{E_L}$ and
Eq.~(\ref{eq:14-gradient}) vanishes identically.  It has to: at an exact
eigenstate the energy is stationary against \emph{any} variation, parameters
included.  Note the practical corollary -- as the trial function improves, the
gradient and the variance vanish together, so a small gradient with a large
variance means the optimiser has found a flat direction rather than the
answer.

\subsection*{Exercise 2: descent methods}
\label{ex:14-descent}
\addcontentsline{toc}{subsection}{Exercise 2: descent methods}

\begin{enumerate}
  \item[a)] For $f(x)=\tfrac12 x^{T}Ax - b^{T}x$ with $A$ symmetric positive
  definite, show that steepest descent with the optimal step converges with a
  rate governed by the condition number of $A$.
  \item[b)] Show that the secant condition~(\ref{eq:14-secant}) is exactly
  satisfied by the true Hessian for a quadratic $f$.
  \item[c)] Why can a quasi-Newton method behave badly when the gradient is
  stochastic?
\end{enumerate}

\paragraph{Answer to a).}
The gradient is $\nabla f = Ax-b$ and the minimum is at $x_{*}=A^{-1}b$.
Writing $e_k=x_k-x_{*}$, one line of algebra gives
$e_{k+1}=(I-\eta A)e_k$, so in the eigenbasis of $A$ each component decays as
$(1-\eta\lambda_i)^k$.  The best uniform $\eta$ balances the extreme
eigenvalues, $\eta=2/(\lambda_{\min}+\lambda_{\max})$, and gives a rate
$(\kappa-1)/(\kappa+1)$ with $\kappa=\lambda_{\max}/\lambda_{\min}$.  A
condition number of $100$ therefore needs of order a hundred times more
iterations than a well-conditioned one -- which is the zig-zag, quantified.

\paragraph{Answer to b).}
For $f$ quadratic with Hessian $A$, $\nabla f(x)=Ax-b$, so
$\nabla f(x_{k+1})-\nabla f(x_k) = A(x_{k+1}-x_k)$, which is
Eq.~(\ref{eq:14-secant}) with $\bm{H}_{k+1}=A$.  Broyden and BFGS impose this
one condition and fix the remaining freedom by asking for the smallest change
to $\bm{H}$ in a suitable norm -- BFGS additionally preserving symmetry and
positive definiteness, which is what makes it suitable for minimisation rather
than root finding.

\paragraph{Answer to c).}
The curvature is inferred from \emph{differences} of gradients, and
differencing two noisy quantities amplifies the noise.  If the statistical
error in $\nabla E$ is comparable to the change in $\nabla E$ between
iterations -- which happens as soon as the steps get small, that is, near the
minimum -- the inferred Hessian is dominated by noise and can easily fail to
be positive definite, sending the optimiser uphill.  The usual remedies are
to increase the number of samples as the optimisation proceeds, to fall back
on a first-order method near the minimum, or to use stochastic
reconfiguration, whose metric~(\ref{eq:14-metric}) is an average rather than a
difference and is therefore far better conditioned statistically.

\subsection*{Exercise 3: resampling}
\label{ex:14-resampling}
\addcontentsline{toc}{subsection}{Exercise 3: resampling}

\begin{enumerate}
  \item[a)] Show that a blocking transformation~(\ref{eq:14-blocking}) leaves
  the sample mean unchanged, and that for uncorrelated data it leaves the
  estimated error unchanged too.
  \item[b)] Show that the jackknife error for the mean is identical to the
  blocked standard error.
  \item[c)] Explain why a bootstrap that resamples individual points rather
  than blocks reproduces the naive error, and what that implies about using it
  on Monte Carlo data.
\end{enumerate}

\paragraph{Answer to a).}
The mean of the blocked data is
$\frac{2}{n}\sum_i\frac12(x_{2i-1}+x_{2i})=\frac1n\sum_j x_j$, unchanged.  For
independent data with variance $\sigma^2$, each block mean has variance
$\sigma^2/2$ and there are $n/2$ of them, so the estimated variance of the
mean is $(\sigma^2/2)/(n/2)=\sigma^2/n$: also unchanged.  The blocking curve
is therefore flat for independent data, and any rise in it is a direct measure
of correlation -- which is what makes the plateau meaningful.

\paragraph{Answer to b).}
With $n$ blocks of means $b_i$ and $\bar b$ their average,
$\hat\theta_{(i)} = (n\bar b - b_i)/(n-1)$, so
$\hat\theta_{(i)}-\bar\theta = -(b_i-\bar b)/(n-1)$.  Substituting into
Eq.~(\ref{eq:14-jackknife}),
\[
  \sigma^2_{\mathrm{jack}}
   = \frac{n-1}{n}\sum_i\frac{(b_i-\bar b)^2}{(n-1)^2}
   = \frac{1}{n(n-1)}\sum_i (b_i-\bar b)^2 ,
\]
which is the ordinary standard error of the $n$ block means.  The jackknife
adds nothing for a linear estimator, and Table~\ref{tab:14-resampling} shows
the two agreeing to three digits.

\paragraph{Answer to c).}
Resampling individual points with replacement produces data sets whose
elements are independent by construction, whatever the original series looked
like.  The bootstrap distribution of the mean then has variance
$\sigma^2/n$ -- the naive answer -- because the correlations have been
destroyed by the resampling itself.  The implication is that a bootstrap
applied naively to Monte Carlo data is not merely inaccurate, it is
systematically and confidently wrong, and looks respectable while being so.
Block bootstrapping, with a block long compared with $\tau$, is the fix.

\subsection*{Exercise 4: numerical explorations}
\addcontentsline{toc}{subsection}{Exercise 4: numerical explorations}

\begin{enumerate}
  \item[a)] Using \texttt{vmcoptimise.py}, reproduce
  Table~\ref{tab:14-descent} and plot the path of the optimiser on top of the
  energy surface of Table~\ref{tab:13-surface}.
  \item[b)] Study the learning rate.  How large may $\eta$ be before the
  descent diverges, and how does that limit relate to the curvature of the
  surface?
  \item[c)] Replace the fixed learning rate by a schedule
  $\eta_k = \eta_0/(1+k/k_0)$ and compare the number of iterations and the
  scatter of the final parameters.
  \item[d)] Feed the analytic gradient to \texttt{scipy.optimize.minimize}
  with \texttt{method="BFGS"} and compare with the hand-written descent.  What
  happens when you reduce the number of cycles per gradient evaluation?
  \item[e)] Verify Exercise~3(c) numerically: bootstrap the production series
  with a block length of one and confirm that you recover the naive error.
  \item[f)] Run the production stage at several time steps and confirm that
  the blocked error, unlike the naive one, is independent of $\Delta t$ -- as
  it must be, since the physical answer cannot depend on the sampling.
\end{enumerate}
